% ===========================================================================
\ReportSection[sec:components]{System Components and Service Expectations}

\ReportSubsection{Plain-language overview}

The system lets many people ask questions of a large language model (LLM) from their
own computers, with the model running entirely on equipment inside the organization.
A person types a question or uploads a document in a web page or a script. A gateway
checks who they are, a retrieval service finds relevant passages in the archived files
on the NAS, and the model, running on a GPU server, writes an answer that cites those
passages. Many people can do this at once because the model software handles their
requests together rather than one after another. A second GPU server is used mainly to
train the model on organization-specific material without slowing anyone down, and it
can also take over serving when needed. The first computer holds the files, backups,
license servers and the supporting software, and nothing leaves the organization's
boundary.

\ReportSubsection{What each piece is and does}

\TableTitle{System components}
{\footnotesize
\begin{xltabular}{\textwidth}{@{}Y{0.8} Y{1.3} Y{1.2} Y{0.7}@{}}
  \toprule
  Component & What it is & What it does for users & Runs on \\
  \midrule
  \endhead
  ML node (GPU server) & A server holding one or more 96~GB graphics processing units (GPUs). The GPU memory is where the model lives while it works. & Provides the computing power to answer questions quickly and to train the model. & Computer~2 \\
  Language model (LLM) & A large file of learned numbers (weights) that predicts text; ``open-weight'' means the organization may run it locally. & Reads questions and documents and writes answers, summaries, code and drafts. & ML node GPUs \\
  Inference server & Software (for example vLLM) that loads the model and serves requests. & Lets many users share one loaded model, streams answers as they are written, and keeps recent context in memory to save time. & ML node \\
  Embedding model and reranker & Small models that turn text into numbers that capture meaning, and that reorder search results by relevance. & Make search work by meaning rather than exact words. & ML node or CPU \\
  Vector database & A database of those numbers for every document passage. & Finds the passages most relevant to a question in milliseconds. & Computer~1 \\
  Ingestion pipeline & Programs that read files on the NAS, extract text (with OCR for scans), split it into passages, and index it. & Make archived documents and standards searchable; re-run when files change. & Node~B or CPU; results stored on Computer~1 \\
  Retrieval service & The logic that takes a question, finds passages, builds the prompt, calls the model, and returns an answer with citations. & Gives answers tied to documents, clauses and pages, and applies each user's access rights. & Computer~1 \\
  LLM gateway & A single front door with an OpenAI-compatible interface. & Signs users in, applies quotas and priorities, routes to Node~A or B, fails over, and logs each request. & Computer~1 \\
  Chat interface & A web page for chatting, uploading files and reusing prompts. & The everyday way most people use the system. & Computer~1 \\
  Training environment & Job queue, fine-tuning tools, experiment tracking and evaluation tests. & Adapts the model to organization vocabulary and formats; tests whether it helped. & Node~B \\
  Model and artifact registry & A versioned, access-controlled store of approved models, adapters and indexes with checksums. & Ensures only approved, verified models run and allows rollback. & Computer~1 \\
  Identity provider & The directory and sign-in service (for example Active Directory with CAC/PIV or multi-factor). & Controls who can use the system and which documents they can retrieve. & Existing or Computer~1 \\
  Logging and monitoring & Collection of audit logs and performance measurements. & Shows who asked what, supports compliance evidence, and warns of overload or failures. & Computer~1, forwarded to the SIEM \\
  File share (NAS) & Shared storage with permissions, snapshots and quotas, served over SMB and NFS. & Holds project data, the archive, datasets, checkpoints and backups. & Computer~1 \\
  License servers & Programs that hand out floating software licenses on the network. & Allow engineering software to run; needed only for network-type licenses. & Computer~1 (VM) \\
  Private cloud (hypervisor) & Virtualization software that runs several isolated virtual machines on one server. & Hosts Autodesk, Bluebeam and other specific software in separate, snapshot-protected machines. & Computer~1 \\
  Storage network and firewall & High-speed switch (25~GbE) and a firewall around the system. & Move data between computers quickly; keep the system inside the boundary. & Both \\
  \bottomrule
\end{xltabular}}

\ReportSubsection{Expected service levels (planning targets)}

These are targets used for sizing and for acceptance testing in the pilot. They are not
guarantees until measured on the real hardware with real prompts.

\TableTitle{Service expectations}
{\footnotesize
\begin{xltabular}{\textwidth}{@{}Y{1.0} Y{0.9} Y{1.1}@{}}
  \toprule
  Measure & Target & Notes \\
  \midrule
  \endhead
  Named users & At least 20 (stated) & Growth: \PH{expected named users in 3 years}. \\
  Simultaneous users generating answers & About 12 (stated) & D3 Node~A is rated 8--20; D4 or Node~B serving adds headroom. \\
  Context size per request & At least 128K tokens (about 300 pages of text) & Quality of recall falls with length; retrieval is preferred for very large archives. \\
  Time to first word & Under 3~s for prompts up to 16K tokens; under 20~s for 128K & Measured at the gateway under load. \\
  Streaming speed per user & At least 15--20 tokens/s (faster than reading speed) & Falls if many long prompts arrive together. \\
  Answers from the archive & Cite document, clause and page & An answer without a citation should be treated as unsupported. \\
  Availability in duty hours & Target \PH{98\%} & Single-node configurations have no ML failover; Computer~1 has no server redundancy. \\
  Recovery after failure & Target \PH{RTO} for LLM serving; \PH{RTO/RPO} for files and licenses & Dual-node: switch to Node~B; test before use and after major changes. \\
  Model refresh & Quarterly review of newer approved models & Tested on Node~B before production. \\
  Fine-tuning turnaround & Hours to a few days per adapter on one GPU & Measure in the pilot; depends on data size. \\
  Archive ingestion rate & \PH{documents per hour, measure in pilot} & OCR of scanned pages is the slowest step. \\
  Data handling & All CUI stays inside the authorization boundary; per-user audit logs & Derived data (indexes, adapters) is treated as CUI. \\
  \bottomrule
\end{xltabular}}

\ReportSubsection{What the system will and will not do}

\textbf{Expected to do well}
\begin{itemize}
  \item Answer questions from approved documents and cite the source passage.
  \item Summarize, compare and extract data from long documents, up to the context limit.
  \item Draft reports, procedures, emails, scripts and code for an engineer to review.
  \item Process large batches of documents overnight.
  \item Learn organization vocabulary, formats and task habits through fine-tuning.
  \item Serve about a dozen people at the same time with acceptable speed.
\end{itemize}

\textbf{Limits to set expectations}
\begin{itemize}
  \item It can be wrong. Every answer used in a deliverable needs engineer review; the
    system does not replace professional judgment or stamping.
  \item Training on a document does not reliably make the model remember it exactly;
    retrieval with citations is how exact wording is obtained.
  \item In the base plan it cannot read drawings or photographs as images; it reads
    text, tables and drawing metadata only.
  \item It does not browse the internet or call outside services.
  \item Pre-training a new model from scratch is out of scope at this budget.
  \item The best open-weight models can trail the best proprietary models on the hardest
    reasoning tasks.
  \item Capacity is finite: at peak, requests queue, and failure of the serving node
    reduces service until failover or repair.
\end{itemize}
